\section{Propagation des ondes}
La propagation des ondes est un phénomène naturel qui peut se produire dans
certains milieux. Dans différentes conditions, l'étude d’une telle propagation
peut engendrer des informations sur la nature des ondes, leurs caractéristiques,
et sur le milieu de propagation.

La figure ci-dessous donne deux dispositifs~(1) et~(2) permettant d'étudier la
propagation d'une onde à la surface de l'eau et la propagation du son dans
l'air.

%\begin{figure}[H]
%    \centering
%    \begin{subfigure}{0.45\textwidth}
%        \centering
%        \includegraphics[width=\textwitdh]{g8136.png}
%        \caption{}
%        \label{fig:}
%    \end{subfigure}
%    \begin{subfigure}{0.45\textwidth}
%        \centering
%        \includegraphics[width=\textwitdh]{g8137.png}
%        \caption{}
%        \label{fig:}
%    \end{subfigure}
%    \caption{}
%    \label{fig:}
%\end{figure}

\begin{figure}[H]
    \centering
    \includegraphics[width=0.7\textwidth]{2025-11-01_012313-}
\end{figure}

\begin{enumerate}
    \item Quelle est la nature de l'onde mécanique produite respectivement par
          les sources de ces deux dispositifs~?

    \item Dans le dispositif~(1), un vibreur produit une onde progressive
          sinusoïdale de fréquence $N_{1}$. Une étude expérimentale a permis
          d'obtenir le document~(a) représentant l'élongation d'un point $M$ de
          la surface de l'eau en fonction du temps et le document~(b)
          représentant l'aspect de la surface de l'eau à un instant donné.
          \begin{figure}[H]
              \centering
              \includegraphics[width=0.8\textwidth]{2025_10_31_756318f153383b902639g-1}
          \end{figure}
          \begin{enumerate}
              \item Lequel des deux documents~(a) et~(b) montre une
                    périodicité spatiale~?
              \item Déterminer la fréquence $N_{1}$ de l'onde.
              \item Calculer la célérité $v_{1}$ de propagation de l'onde à la
                    surface de l'eau.
              \item Recopier sur votre copie le numéro de la question et écrire
                    la lettre correspondante à la proposition vraie.
              
                    L'élongation du point $M$ en fonction de l'élongation de la
                    source $S$ s'écrit~:
                    \begin{table}[H]
                        \flushright
                        \setlength{\arrayrulewidth}{1pt}
                        \renewcommand{\arraystretch}{1.5}
                        \begin{tabularx}{0.95\textwidth}{|p{3mm}|X|p{3mm}|X|p{3mm}|X|p{3mm}|X|}
                            \hline
                            $A$ & $y_{M}(t)=y_{S}(t+0,1)$ &
                            $B$ & $y_{M}(t)=y_{S}(t+0,05)$ &
                            $C$ & $y_{M}(t)=y_{S}(t-0,1)$ &
                            $D$ & $y_{M}(t)=y_{S}(t-0,05)$ \\
                            \hline
                        \end{tabularx}
                    \end{table}
          \end{enumerate}
    \item On interpose à la surface de l'eau un obstacle muni d'une ouverture de
          diamètre $L=\qty{8}{\cm}$. L'onde produite à la surface de l'eau par
          la source se propage après avoir traversé l'ouverture.
          \begin{enumerate}
              \item Quel phénomène peut-on observer lorsque l'onde traverse
                    l'ouverture~? Justifier.
              \item Déduire la longueur d'onde $\lambda_{2}$ et la célérité de
                    propagation $v_{2}$ de l'onde au-delà de l'ouverture.
          \end{enumerate}
    \item Le haut-parleur du dispositif~(2), émet des ondes sonores de fréquence
          $N_{2}=\qty{10}{\kHz}$.
          \begin{enumerate}
              \item Les ondes sonores produites peuvent-elles se propager dans
                    le vide? Justifier.
              \item Les ondes sont captées par deux microphones $M_{1}$
                    et $M_{2}$ qui occupent la même position. Les
                    courbes visualisées sur l'écran de l'oscilloscope
                    apparaissent en phase.
              
                    Lorsqu'on déplace $M_{2}$ par rapport à
                    $M_{1}$ d'une distance $d=\qty{34}{\cm}$, les deux courbes
                    observées à l'oscilloscope apparaissent à nouveau en phase
                    pour la dixième~(10) fois. Déduire la célérité de
                    propagation du son dans l'air.
          \end{enumerate}
\end{enumerate}



\section*{Exercice 2 ( 2,5 points) : Transformations nucléaires}
La radioactivité est un phénomène naturel et durable produit par des sources radioactives. Suite à des désintégrations en chaine, un nucléide peut se transformer en d'autres jusqu'à obtention d'un nucléide stable, formant ainsi une famille radioactive. Selon leurs durées de vie, ces sources peuvent avoir des avantages et des inconvénients.

Le diagramme ci-contre donne quelques nucléides appartenant à la famille radioactive de l'uranium.\\
\includegraphics[width=\textwidth]{2025_11_26_673124dc142c0df571a8g-4(2)}

\section*{Données :}
$\mathrm{m}\left({ }_{83}^{212} \mathrm{Bi}\right)=211,94562 u ; \mathrm{m}\left({ }_{81}^{208} \mathrm{Tl}\right)=207,93745 u ; \mathrm{m}(\alpha)=4,00150 u ; 1 u=931,5 \mathrm{MeV} . \mathrm{c}^{-2}$

\begin{enumerate}
  \item Préciser, en justifiant, si les nucléides ${ }_{82}^{212} \mathrm{~Pb}$ et ${ }_{83}^{212} \mathrm{Bi}$ sont des isotopes.
  \item Identifier, en justifiant, le type de la désintégration (1) (voir diagramme).
  \item Reconnaitre le nucléide ${ }_{Z}^{A} X$.
  \item Déterminer, en unité ( $M e V$ ), la valeur de l'énergie libérée $E_{\text {libérée }}=|\Delta E|$ par la désintégration d'un noyau de bismuth ${ }_{83}^{212} B i$ en thallium ${ }_{81}^{208} T l$.
  \item Soit une source radioactive contenant à l'instant $\left(t_{0}=0\right), \mathrm{N}_{0}=28,4.10^{19}$ noyaux de bismuth ${ }_{83}^{212} \mathrm{Bi}$ radioactif. Pendant la durée de 15 min , un compteur a enregistré $4,484.10^{19}$ désintégrations.\\
5.1. Quelle est le nombre de noyaux de bismuth ${ }_{83}^{212} \mathrm{Bi}$ présent dans la source à l'instant $\mathrm{t}_{1}=15 \mathrm{~min}$ ?\\
5.2. Déterminer la période radioactive (demi-vie) $\mathrm{t}_{1 / 2}$ du bismuth ${ }_{83}^{212} B i$.\\
5.3. Le nucléide de bismuth ${ }_{83}^{212} B i$ peut-t-il être utilisé pour la datation d'un événement? Justifier.
\end{enumerate}

\section*{Exercice 3 (6,5 points): Dipôle RC - Circuit RLC série}
Les condensateurs sont des composants électroniques qu'on retrouve dans plusieurs circuits électriques et électroniques. Ils diffèrent par leurs formes et leurs technologies. Placés dans des circuits, ils permettent un stockage de l'énergie. Cette énergie est d'autant plus importante pour des condensateurs de grande capacité, et qui sera transférée lors de différents usages et utilisations.

Le montage de la figure 1 comporte :

\begin{itemize}
  \item un générateur idéal de tension de force électromotrice E;
  \item un condensateur de capacité C réglable;
  \item un conducteur ohmique de résistance $R$;
  \item une bobine d'inductance $L$ et de résistance $r$;
  \item un interrupteur $K$ à double position.
\end{itemize}

Données : $R=100 \Omega ; r=20 \Omega$

\section*{Partie 1: Étude de la charge du condensateur}
À l'instant $t_{0}=0$, on place l'interrupteur $K$ en position 1 .

\begin{figure}[h]
\begin{center}
  \includegraphics[width=\textwidth]{2025_11_26_673124dc142c0df571a8g-5}
\captionsetup{labelformat=empty}
\caption{Figure 1}
\end{center}
\end{figure}

\begin{enumerate}
  \item Préciser l'intérêt du montage représenté sur la figure 1 (l'interrupteur K est en position 1 ).
  \item Établir l'équation différentielle vérifiée par la tension $u_{C}(t)$ aux bornes du condensateur.
  \item À l'aide d'un système d'acquisition convenable, on obtient les courbes (1) et (2) de la figure 2 qui représentent l'évolution temporelle de la tension $u_{C}(t)$ pour deux valeurs $C_{1}$ et $C_{2}$ de la capacité. On désigne par $\tau_{1}$ et $\tau_{2}$ respectivement les constantes de temps relatives aux courbes (1) et (2).\\
3.1. Déterminer pour la courbe (2), la durée du régime transitoire.\\
3.2. Calculer les valeurs de $C_{1}$ et $C_{2}$.\\
3.3. Quelle est l'influence de la
\end{enumerate}

\begin{figure}[h]
\begin{center}
  \includegraphics[width=\textwidth]{2025_11_26_673124dc142c0df571a8g-5(1)}
\captionsetup{labelformat=empty}
\caption{Figure 2}
\end{center}
\end{figure}

3.4. Déterminer la valeur de la force électromotrice $E$.\\
3.5. Déterminer, pour le condensateur de capacité $C_{1}$, la valeur de sa charge $q_{1}$ à l'instant $t=\tau_{1}$.\\
3.6. En utilisant le même générateur de force électromotrice $E$, préciser dans quel cas ( $C_{1}$ ou $C_{2}$ ), le condensateur emmagasinera la plus grande énergie électrique à la fin de la charge. Justifier.

\section*{Partie 2 : Étude d'un circuit RLC série}
Le condensateur réglé sur la capacité $C=1 \mu F$ est chargé complètement. On réalise sa décharge à travers la bobine en basculant $K$ en position 2 . À l'aide du même système d'acquisition, on obtient la courbe de la figure 3 qui représente $u_{C}(t)$.

\begin{figure}[h]
\begin{center}
  \includegraphics[width=\textwidth]{2025_11_26_673124dc142c0df571a8g-6}
\captionsetup{labelformat=empty}
\caption{Figure 3}
\end{center}
\end{figure}

\begin{enumerate}
  \item Expliquer qualitativement la variation de l'amplitude des oscillations.
  \item Quelle est la valeur de la pseudo- période $T$ des oscillations ?
  \item Déduire la valeur de l'inductance $L$ de la bobine sachant que la pseudo- période est égale à la période propre de l'oscillateur (LC).
  \item Pour entretenir les oscillations électriques dans le circuit RLC série, on monte en série avec le condensateur et la bobine, un générateur G délivrant une tension $u_{g}$ proportionnelle à l'intensité $i(t)$ du courant électrique ( $u_{g}=k . i$ ). ( $k$ constante positive).\\
4.1. Quel est le rôle du générateur G de point de vue énergétique ?\\
4.2. Quelle valeur doit prendre $k$ pour obtenir des oscillations entretenues ? Justifier.\\
4.3. Que peut-on dire des oscillations électriques obtenues après l'entretien?
\end{enumerate}



